% ============================================================
%  Figure captions for the four-route convergence paper.
%  Panels 5-8: the ladder and convergence routes.
%  Every number quoted here is read from results/ladder_routes.json.
% ============================================================

\newcommand{\captionFourRoutes}{%
\caption{\textbf{Four Routes to One Value, and What the Cross-Line
Constraint Catches.}
The four routes share the alphabet $C(n)=2n^{2}$ and nothing else: the
established route solves a differential equation, the instrument route
resolves a partition address physically, the ladder route deletes the
carrier and computes invariants of a contact sequence, and the catalogue
route takes a minimum cut at rest.
(\textbf{A})~Signed residual from the NIST measurement, in parts per
million, for each route on each Lyman line. The established, instrument and
catalogue routes are \emph{bit-identical} at $-11$~ppm; the ladder bar is
zero because it is reconstructed from the measured value. Plotting the raw
wavenumbers instead would show four bars of equal height and say nothing,
which is why the residual is drawn.
(\textbf{B})~The spread \emph{among} the three predictive routes against
their common gap to measurement. The spread is exactly zero and is drawn at
a clamped floor with that fact labelled, because rendering machine zero on a
logarithmic axis would dress floating-point dust as data. The honest reading
of (\textbf{A}) and (\textbf{B}) together is that the three routes are three
\emph{derivations of one closed form}, not three independent predictions
that happen to agree; the shared $1.1\times10^{-5}$ residual is the
reduced-mass and QED gap between that closed form and experiment.
(\textbf{C})~The measurement that a single route cannot make. Ritz
additivity, $\tilde\nu(1\!\to\!n)=\tilde\nu(1\!\to\!2)+\tilde\nu(2\!\to\!n)$,
evaluated on the tabulated wavelengths (red) and after converting the Balmer
values from standard air to vacuum (green), against the rounding budget of
the quoted four-decimal digits (dotted). As tabulated the residual is
$6.4\times10^{-5}$, exceeding the budget by a factor of $73$ and rising
monotonically with $n$; converted, it falls to $7.1\times10^{-7}$, inside the
budget, and the trend disappears. The tabulated Balmer wavelengths are
standard-air values sitting alongside vacuum Lyman values. A per-line
comparison cannot detect this, because every line agrees with the table it
was drawn from; the cross-line constraint detects it immediately.
(\textbf{D})~Residual landscape over (route, line, $\log_{10}$ relative
residual), showing the three coincident routes as three parallel planes and
the reconstructed ladder route on the floor.}}

\newcommand{\captionLadderRungs}{%
\caption{\textbf{Rung Powers and Circulation Computed from Measured
Lines.}
A transition is a rung; its power $\pi$ is the fraction of the final
state's binding that the transition supplies, computed throughout from the
measured NIST wavenumber rather than assigned. This is the half the
categorical-ladder development did not have: its ring profiles were inputs,
so its chemical claims were statements about a formalism.
(\textbf{A})~Rung power against upper level for the five hydrogen series.
Each series is an ordered fan approaching $\pi\to1$ at its limit, where the
transition supplies the whole of the final state's binding.
(\textbf{B})~The circulation contribution $-\log(1-\pi)$ of each rung. The
divergence toward the series limit is the statement that a rung closing all
remaining ambiguity deposits unbounded residue.
(\textbf{C})~Total circulation $\varrho$ and circulation per rung
$\varrho/n$ by series. $\varrho$ is additive along the ladder, which is the
ladder-language form of the Ritz combination principle tested in
Figure~\ref{fig:fourroutes}(C).
(\textbf{D})~Rung landscape over $(n_f, n_u, \pi)$, the whole measured
hydrogen atlas as a single surface in power coordinates.}}

\newcommand{\captionHyperfineRung}{%
\caption{\textbf{The 21\,cm Line as the Parity Rung: Nearly Inert, and
Therefore Nearly Forbidden.}
A rung rewrites one or more letters of the word $(n,\ell,m,s)$. The
hyperfine transition rewrites the parity letter $s$ alone, with $n$, $\ell$
and $m$ all fixed.
(\textbf{A})~Every allowed electric-dipole rung of the measured hydrogen
atlas, sorted by power, against the 21\,cm rung and against the inert
(forbidden) line at exactly zero. The parity rung sits five decades below
the body of the atlas and above zero.
(\textbf{B})~The three cases side by side on a logarithmic power axis: a
forbidden rung at exactly $0$, the 21\,cm rung at
$\pi=4.32\times10^{-7}$, and the smallest allowed rung at $\pi=0.306$. The
parity rung lies $5.85$ orders of magnitude below the smallest allowed one.
The registered expectation was ``at least four orders''; the measured value
exceeds it.
(\textbf{C})~Which letter each rung rewrites. The two allowed transitions
move $n$ and $\ell$ together, satisfying $\Delta\ell=\pm1$; the
$2s\to1s$ transition moves $n$ alone and is dipole-forbidden; the 21\,cm
rung moves $s$ alone. Since chirality is a topological invariant and
continuous deformation cannot change it, a rung on $s$ alone is forbidden as
an electric-dipole transition --- and the near-vanishing power in
(\textbf{B}) is the quantitative form of that statement, and of the
$10^{7}$-year lifetime that follows from it.
(\textbf{D})~Ten decades of transition energy in
$(\log_{10}\tilde\nu, n_f, \log_{10}\pi)$, with the parity rung starred at
the far corner.}}

\newcommand{\captionClosedLadders}{%
\caption{\textbf{Closed Molecular Ladders, with Powers from HITRAN.}
A cycle has no target, so composite power is undefined for it; the closed
ladder carries the circulation $\varrho$ and the uniformity $u$ instead.
(\textbf{A})~The three vibrational modes of H$_2$O drawn as a closed
ladder, node area and colour encoding the rung power computed from the
measured HITRAN fundamentals against the dissociation limit. The bend
$\nu_2$ at $0.0388$ is the odd node against the two stretches at $0.0889$
and $0.0913$, and that asymmetry is exactly what gives $u=0.6682$.
(\textbf{B})~Both invariants under every rotation of the cycle. Each is
flat, with deviation below $10^{-16}$. The values are plotted rather than
the deviations, because a deviation of machine zero on a logarithmic axis
would render floating-point dust as a measurement.
(\textbf{C})~Circulation per rung for the electronic (Lyman) and
vibrational (H$_2$O) modalities, separated with margin $+2.556$. The
corresponding separation in the ring corpus of the categorical-ladder
development was $+0.293$ on assigned profiles; here the profiles are
measured.
(\textbf{D})~The closed-ladder coordinates $(\varrho/n, u, \text{cycle
length})$ for the four cycles treated: the H$_2$O vibrational triple, the
H$_2$O rotational constants $(A,B,C)$, the single H$_2$ vibrational mode,
and the Lyman ladder. Neither coordinate alone separates the four --- the
same two-invariant conclusion the ring corpus reached when uniformity alone
was refuted by cyclohexane.}}
